\documentclass[10pt,a4paper]{amsart}
\usepackage[margin=19mm]{geometry}
\usepackage{amsmath,amssymb,mathtools}
\usepackage[scr=rsfs,cal=euler]{mathalfa}
\usepackage{xfrac}
\usepackage[unicode,hidelinks,bookmarks=false]{hyperref}
\newcommand{\ie}{i.e.\ }
\newcommand{\cf}{cf.\ }
\newcommand{\resp}{respectively\ }
\newcommand{\Subsection}{\subsection}
\providecommand{\nobreakdash}{\nobreak\hbox{-}}
\setlength{\emergencystretch}{2em}
\multlinegap0pt
\usepackage{bm}
\usepackage{bbm}
\usepackage{dsfont}
\usepackage{xspace}
\usepackage{cancel}
\usepackage{mathtools}
\usepackage{amsthm}
\makeatletter
\@ifundefined{theorem}{\newtheorem{theorem}{Theorem}}{}
\@ifundefined{lemma}{\newtheorem{lemma}[theorem]{Lemma}}{}
\makeatother
\newcommand{\dist}{\operatorname{dist}}
\title[Lipschitz regularity of harmonic quasiconformal maps between]{Lipschitz regularity of harmonic quasiconformal maps between Lyapunov domains in $\mathbb{R}^n$}
\author{Anton Gjokaj, David Kalaj}
\date{Source: arXiv:2602.05436v1 (CC BY 4.0)}
\begin{document}
\maketitle

\noindent\emph{Source and licence.} Lipschitz regularity of harmonic quasiconformal maps between Lyapunov domains in $\mathbb{R}^n$. Version arXiv:2602.05436v1 (CC BY 4.0), distributed under the Creative Commons Attribution 4.0 International licence, as recorded in the arXiv licence metadata for this exact version and shown on the abstract page https://arxiv.org/abs/2602.05436v1. Full rights notice: licenses/arxiv-2602.05436-CC-BY-4.0.txt.\par\smallskip

\noindent\textit{Editorial scope.} Counted unit: the Lemma whose source label is \texttt{lema2} in the article (source0.tex lines 380--407, source label \texttt{lema2}), delivered together with the article's own complete original proof of it (source0.tex lines 408--429, one \texttt{proof} environment of 22 lines). No mathematics is written, repaired or re-derived: statement and proof are the article's own text line for line. Required context retained uncounted: lem:anchored\_holder\_to\_grad\_normal (lines 227--270); eq:anchored\_holder\_bd\_combined (lines 238--240); eq:Fy\_decay\_combined (lines 333--335). Every definition, display and label of those lines is retained unchanged. Omitted from this package and not redistributed: 1--226, 241--332, 336--379, 430--922. Nothing inside a retained excerpt is omitted or altered: every retained body is delivered exactly as the article's lines are written, so no omission receipt of any kind accompanies this package. Citations: the retained excerpts carry no citation command at all, so no bibliography entry of the article is transcribed and no reference is redistributed. Reference adaptation: none was needed and none was made; every reference inside the delivered unit resolves inside it, so no unresolved reference or citation is produced. Printed slips kept literal: no printed word, symbol, hypothesis or number of the counted statement or of its proof was altered. Layout and class: the article's own class is \texttt{amsart} with packages including amsmath, amssymb, amsthm, graphicx, titlesec, fancyhdr; the wrapper is the lane's pinned \texttt{amsart} editorial wrapper at 10pt on A4, which restates the theorem-like environments the retained text needs and adds the article's own macro definitions that the retained text needs (\texttt{\textbackslash dist}), with extra wrapper packages (bm, bbm, dsfont, xspace, cancel, mathtools). The delivered numbering is the unit's own. Assets: no retained span contains a figure environment, a graphics-inclusion command, a PDF-page inclusion, a table or a TikZ picture; no image file is copied into this package, the package declares no asset dependency and no image file is redistributed with it. Licence: version arXiv:2602.05436v1 of this article is distributed under the Creative Commons Attribution 4.0 International licence, as recorded in the arXiv licence metadata for this exact version and shown on the abstract page https://arxiv.org/abs/2602.05436v1. A full rights notice is delivered at licenses/arxiv-2602.05436-CC-BY-4.0.txt. Characters: every retained excerpt is pure ASCII, so the delivered engine, XeLaTeX with its default Latin Modern text font, reports no missing character.
\begin{lemma}[Basepoint boundary H\"older at $x_0$ $\Rightarrow$ interior H\"older and gradient control along the normal]
\label{lem:anchored_holder_to_grad_normal}
Let $D\subset\mathbf R^n$ be a domain and let $x_0\in\partial D$.
Assume that $\partial D$ is $C^{1,\alpha}$ in a neighborhood of $x_0$ for some $\alpha\in(0,1)$.
Fix $0<\mu< 1$.

Then there exist $r_0>0$, $c\in(0,1)$ and $C>0$ (depending only on the local $C^{1,\alpha}$ character
of $\partial D$ near $x_0$ and on $\mu$) such that the following holds.

Let $u$ be harmonic in $\Omega_{x_0}:=D\cap B(x_0,r_0)$ and continuous on $\overline{\Omega_{x_0}}$.
Assume that on the boundary patch $\Gamma:=\partial D\cap B(x_0,r_0)$,
\begin{equation}\label{eq:anchored_holder_bd_combined}
|u(\xi)-u(x_0)|\le M|\xi-x_0|^\mu,\qquad \xi\in\Gamma,
\end{equation}
and set $A:=\|u\|_{L^\infty(\Omega_{x_0})}$.
Let $\nu(x_0)$ be the inward unit normal to $\partial D$ at $x_0$.
\\
Then for every
\[
x=x_0+t\,\nu(x_0)\in D,\qquad 0<t<cr_0,
\]
we have the ``tent'' estimate
\begin{equation}\label{eq:tent_holder_combined}
\sup_{y\in D\cap B(x,t/2)} |u(y)-u(x_0)|
\;\le\; C\Big(M+\frac{A}{r_0^\mu}\Big)\,t^\mu.
\end{equation}
In particular,
\begin{equation}\label{eq:holder_normal_combined}
|u(x)-u(x_0)|
\;\le\; C\Big(M+\frac{A}{r_0^\mu}\Big)\,t^\mu,
\end{equation}
and the full gradient satisfies
\begin{equation}\label{eq:grad_normal_combined}
|\nabla u(x)|
\;\le\; C\Big(M+\frac{A}{r_0^\mu}\Big)\,t^{\mu-1}.
\end{equation}
Equivalently, writing $\delta_D(x):=\dist(x,\partial D)$, and noting that
$\delta_D(x)=t$ along this inward normal ray for $t<cr_0$, we obtain
\[
|u(x)-u(x_0)|\le C\Big(M+\frac{A}{r_0^\mu}\Big)\,\delta_D(x)^\mu,
\qquad
|\nabla u(x)|\le C\Big(M+\frac{A}{r_0^\mu}\Big)\,\delta_D(x)^{\mu-1}.
\]
\end{lemma}

\begin{equation}\label{eq:Fy_decay_combined}
F_y(s)\le C\,\frac{\delta_D(y)}{s},
\end{equation}

\begin{lemma}[Addendum to Lemma~\ref{lem:anchored_holder_to_grad_normal}: the case $\mu>1$]
\label{lema2}
Assume the hypotheses and notation of Lemma~\ref{lem:anchored_holder_to_grad_normal}, except that
\eqref{eq:anchored_holder_bd_combined} holds for some exponent $\mu>1$.
Then there exist $r_0>0$, $c\in(0,1)$ and $C>0$ (depending only on the local $C^{1,\alpha}$ character
of $\partial D$ near $x_0$ and on $\mu$) such that for every
\[
x=x_0+t\,\nu(x_0)\in D,\qquad 0<t<cr_0,
\]
we have
\begin{equation}\label{eq:holder_mu_gt_1}
|u(x)-u(x_0)|
\;\le\; C\Big(M\,r_0^{\mu-1}+\frac{A}{r_0}\Big)\,t,
\end{equation}
and
\begin{equation}\label{eq:grad_mu_gt_1}
|\nabla u(x)|
\;\le\; C\Big(M\,r_0^{\mu-1}+\frac{A}{r_0}\Big).
\end{equation}
Equivalently, along this inward normal ray (where $\delta_D(x)=t$ for $t<cr_0$),
\[
|u(x)-u(x_0)|
\le C\Big(M\,r_0^{\mu-1}+\frac{A}{r_0}\Big)\,\delta_D(x),
\qquad
|\nabla u(x)|
\le C\Big(M\,r_0^{\mu-1}+\frac{A}{r_0}\Big).
\]
\end{lemma}

\begin{proof}[Proof]
The proof is identical to that of Lemma~\ref{lem:anchored_holder_to_grad_normal}, except for the estimate of
\[
\int_{\Gamma}|\zeta|^\mu\,d\omega^y(\zeta).
\]
Using the same notation and the layer--cake representation,
\[
\int_{\Gamma}|\zeta|^\mu\,d\omega^y(\zeta)=\mu\int_0^1 s^{\mu-1}F_y(s)\,ds,
\]
and the same decay estimate \eqref{eq:Fy_decay_combined} for $s\in(\delta_D(y),1]$,
we split at $s=\delta_D(y)$ to obtain
\begin{align*}
\int_{\Gamma}|\zeta|^\mu\,d\omega^y(\zeta)
&\le \mu\int_0^{\delta_D(y)} s^{\mu-1}\,ds
   + C\mu\,\delta_D(y)\int_{\delta_D(y)}^{1} s^{\mu-2}\,ds \\
&= \delta_D(y)^\mu
   + C\mu\,\delta_D(y)\,\frac{1-\delta_D(y)^{\mu-1}}{\mu-1}
\;\le\; C\,\delta_D(y),
\end{align*}
since $\mu>1$ and $\delta_D(y)\in(0,1)$. Substituting this bound in place of the corresponding estimate
in Lemma~\ref{lem:anchored_holder_to_grad_normal} yields the stated conclusions.
\end{proof}

\end{document}
